\Needspace{12\baselineskip}
\section{References}

\begin{multicols}{2}
\fontsize{8.25}{9.6}\selectfont
\RaggedRight
\setlength{\emergencystretch}{1em}

\subsection{Sacred Scripture and doctrine}

\begin{itemize}
\item Sacred Scripture: Gen 1:1--2; 11:1--9; Ex 19; Lev 23:15--21; Deut 16:9--12; Ps 103(104):30; Isa 6:1--8; 11:1--5; 61:1--3; Jer 31:31--34; Ezek 36:25--27; 37:1--14; Joel 3:1--5; Lk 1:26--38; 24:44--53; Jn 7:37--39; 14--16; 20:19--23; Acts 1--2; 4:23--31; 8; 10; 13; 15--16; Rom 5; 8; 12; 1 Cor 2; 6; 12--14; Gal 4--5; Eph 1; 4; 1 Thess 5:19--22; 1 Pet 4:7--11; 1 Jn 4:1--3.
\item First Council of Constantinople (381), Nicene-Constantinopolitan Creed; Council of Florence, \textit{Laetentur caeli} (1439), on the procession of the Holy Spirit.
\item Second Vatican Council, \textit{Sacrosanctum Concilium} 10--13; \textit{Lumen gentium} 4, 12, 48--49, 60--62; \textit{Ad gentes} 4; \textit{Dei Verbum} 7--10.
\item \textit{Catechism of the Catholic Church}, nos.~243--248, 683--747, 797--801, 1076, 1285--1321, 1695, 1830--1832, 1996--2005, 2623, 2670--2672.
\item St.~John Paul II, \textit{Dominum et vivificantem} (18 May 1986), especially nos.~3--10, 22--26, 51--54, 58--66.
\end{itemize}

\subsection{Liturgy and discipline}

\begin{itemize}
\item Leo XIII, \textit{Provida Matris} (5 May 1895), official Holy See witness; and \textit{Divinum illud munus} (9 May 1897), especially nos.~1--7 and 11--14, \sourceurl{https://www.vatican.va/content/leo-xiii/en/encyclicals/documents/hf_l-xiii_enc_09051897_divinum-illud-munus.html}{official English text}.
\item Congregation for Divine Worship and the Discipline of the Sacraments, \textit{Directory on Popular Piety and the Liturgy} (17 December 2001), nos.~32, 61--64, 76--80, 94, 155--156, \sourceurl{https://www.vatican.va/roman_curia/congregations/ccdds/documents/rc_con_ccdds_doc_20020513_vers-direttorio_en.html}{official English text}.
\item Paul VI, \textit{Mysterii Paschalis} (14 February 1969) and the General Norms for the Liturgical Year and the Calendar, nos.~22--26, for Easter Time, Ascension, and Pentecost.
\item Apostolic Penitentiary, fourth \textit{Enchiridion Indulgentiarum}, general norms and concessions 22 and 26, \sourceurl{https://www.vatican.va/roman_curia/tribunals/apost_penit/documents/rc_trib_appen_doc_20020826_enchiridion-indulgentiarum_lt.html}{official Latin text}.
\item \textit{Compendium of the Catechism of the Catholic Church} (2005), appendix of common prayers, for the received Latin \textit{Veni Creator Spiritus} and \textit{Veni Sancte Spiritus}, \sourceurl{https://www.vatican.va/archive/compendium_ccc/documents/archive_2005_compendium-ccc_en.html}{official bilingual witness}.
\item \textit{Missale Romanum}, editio typica (Vatican City: Typis Polyglottis Vaticanis, 1962): Good Friday order no.~960, p.~180 (facsimile PDF p.~261), for the complete Lord's Prayer; \textit{Ordo Missae}, p.~216 (facsimile PDF p.~297), for the Lesser Doxology; \textit{Rubricae generales} 115(a)--(e), p.~XIX (facsimile PDF p.~17), for the exact full conclusion formulas; Easter Thursday, p.~336; Sunday after Ascension, p.~353; Pentecost Tuesday and Ember Wednesday, pp.~360--361; Ember Friday, p.~365; Ember Saturday's shorter form, p.~370; Annunciation, p.~495; \textit{Orationes diversae} 29, p.~[113]; and \textit{Pro fidei propagatione}, p.~[86], \sourceurl{https://media.churchmusicassociation.org/pdf/missale62.pdf}{typical-edition facsimile}.
\item Dom Gaspar Lefebvre, O.S.B., \textit{Daily Missal with Vespers for Sundays and Feasts} (Lophem near Bruges: St.~Andrew's Abbey; St.~Paul, Minnesota: The E.~M. Lohmann Co., 1925), \sourceurl{https://archive.ccwatershed.org/media/pdfs/21/02/16/05-35-23_0.pdf}{edition scan}: common prayers, pp.~4--5 and 35; conclusion formulas and responses, pp.~8--9, 833, and 861; daily collects, pp.~180--181, 895, 948, 977, 980, 987, 992, 1310, and 1827; daily invocation, p.~1873. This is the governing English prayer witness and the governing Latin witness for the Sign of the Cross on p.~35. Its title matter records Theodorus's abbatial \textit{Imprimatur} (21 March 1924), William Busch's \textit{Nihil Obstat}, and Archbishop Austin Dowling's \textit{Imprimatur} (13 May 1925); those historical permissions do not approve this devotional arrangement or establish a current liturgical translation.
\end{itemize}

\subsection{Patristic and scholastic witnesses}

\begin{itemize}
\item St.~Basil of Caesarea, \textit{On the Holy Spirit}, 9.22--23; 16.37--40; 18.45--47; 26.61--64.
\item St.~Cyril of Jerusalem, \textit{Catechetical Lectures}, 16.11--16 and 17.14--19, on living water, diverse gifts, and Pentecost.
\item St.~John Chrysostom, \textit{Homilies on the Acts of the Apostles}, 3--5, on Acts 1--2.
\item St.~Augustine, \textit{Sermons} 267 and 269, on the one Spirit, many tongues, and the catholic Church; \textit{On the Trinity} IV.20.29, on the Spirit's mission.
\item St.~Leo the Great, \textit{Sermon} 75 on Pentecost and \textit{Sermons} 73--74 on the Ascension.
\item St.~Thomas Aquinas, \textit{Summa theologiae} I, qq.~36 and 43; I--II, qq.~68--70; II--II, q.~1; III, q.~57; and commentary on John 14--16.
\end{itemize}

\subsection{Hymn texts and history}

\begin{itemize}
\item Edward Caswall, \textit{Lyra Catholica: Containing All the Breviary and Missal Hymns, with Others from Various Sources} (London: James Burns, 1849), pp.~103--104, “Come, O Creator Spirit blest!”, and pp.~234--236, “Holy Spirit! Lord of light!”, \sourceurl{https://archive.org/details/LyraCatholica1849}{public-domain edition scan}; governing witness for the complete English verse texts.
\item F.~J.~E. Raby, \textit{A History of Christian-Latin Poetry from the Beginnings to the Close of the Middle Ages}, 2nd ed. (Oxford: Clarendon Press, 1953), on the ninth-century setting and uncertain authorship of \textit{Veni Creator}.
\item John Julian, ed., \textit{A Dictionary of Hymnology}, rev. ed. (London: John Murray, 1907), entries on \textit{Veni Creator Spiritus} and \textit{Veni Sancte Spiritus}; used with later manuscript-catalog cautions for authorship claims.
\item Cantus Index and the Digital Image Archive of Medieval Music, manuscript records for the Pentecost sequence; used to bound rather than settle competing authorship traditions.
\end{itemize}

\end{multicols}
